% ==========================================================
% 平均分配中的概率与选项判断
% 难度：★★ 中档
% ==========================================================

\newcommand{\qProbPermAvgAllocStem}{%
将 $5$ 名男大学生、$4$ 名女大学生平均分配到甲、乙、丙三所学校，每所学校 $3$ 人，判断下列概率结论是否正确。\\
A. 甲学校没有女大学生的概率为 $\frac{5}{21}$\\
B. 甲学校至少有两名女大学生的概率为 $\frac{25}{42}$\\
C. 每所学校都有男大学生的概率为 $\frac{6}{7}$\\
D. 乙学校分配两名女生、一名男生，同时丙学校至少有一名女生的概率为 $\frac{1}{7}$
}

\newcommand{\qProbPermAvgAllocAnswer}{%
C
}

\newcommand{\qProbPermAvgAllocSolution}{%
九名学生互不相同，三所学校有标号。先为甲校选择 $3$ 人，再为乙校选择 $3$ 人，其余进入丙校，因此统一样本空间大小为 $\lvert\Omega\rvert = \binom{9}{3}\binom{6}{3} = 1680$。

\textbf{判断选项 A：} 甲学校没有女大学生，意味着甲校的三人全是男生。有利分配数为 $\binom{5}{3}\binom{6}{3}$。所以 $P = \frac{\binom{5}{3}\binom{6}{3}}{\binom{9}{3}\binom{6}{3}} = \frac{5}{42}$。选项给出的是 $\frac{5}{21}$，A 错误。

\textbf{判断选项 B：} 甲学校至少有两名女大学生，分为甲恰有两名女生和甲恰有三名女生。
有利分配数为 $\left(\binom{4}{2}\binom{5}{1}+\binom{4}{3}\right)\binom{6}{3}$。所以 $P = \frac{680}{1680} = \frac{17}{42}$。选项给出的是 $\frac{25}{42}$，B 错误。

\textbf{判断选项 C：} 直接分类较繁琐，使用补集更简洁。若某一所学校没有男生，则该校的三人全部是女生。由于一共只有四名女生，不可能有两所学校同时没有男生，各坏事件互斥。坏情况数为 $3\binom{4}{3}\binom{6}{3}$。所以 $P = 1 - \frac{3\binom{4}{3}\binom{6}{3}}{1680} = 1 - \frac{1}{7} = \frac{6}{7}$。C 正确。

\textbf{判断选项 D：} 先安排乙校：$\binom{4}{2}\binom{5}{1}$。此时剩余六人中有两名女生、四名男生。丙校选择三人且至少有一名女生，方法数为 $\binom{6}{3}-\binom{4}{3}$。所以 $P = \frac{\binom{4}{2}\binom{5}{1}\left(\binom{6}{3}-\binom{4}{3}\right)}{1680} = \frac{480}{1680} = \frac{2}{7}$。选项给出的是 $\frac{1}{7}$，D 错误。

综上，正确选项为 C。
}

\newcommand{\qProbPermAvgAllocTeachingNotes}{%
\textbf{【避坑与边界说明】}
\begin{itemize}
    \item 甲、乙、丙三所学校不同，不能把三组视为无标号组。
    \item 计算某校条件时，安排好该校后，剩余人员仍需分配到另外两所学校。
    \item “至少”类问题要检查直接分类和补集哪一种更简洁。
    \item 使用补集时，要判断多个坏事件能否同时发生。
\end{itemize}
}
